In this appendix we recall some useful material on compact rigid generation.
Most of this material has appeared elsewhere and all of it is certainly known to experts.
This appendix was mostly included for the convenience of the authors,
though we do hope the reader finds it to be a concise summary of a basic technique in higher algebra.

\begin{dfn}
  A stable, presentably monoidal category $\CC$ is \emph{rigidly generated} if
  it has a family of compact dualizable generators. \footnote{Here and throughout this appendix we only require that an object has a one-sided dual. The side on which an object has a dual will not be important.}
\end{dfn}

Our goal will be to study some properties of monoidal left adjoints out of a rigidly generated category.
This begins with the following construction.

\begin{cnstr}
  Given an $\En$-monoidal left adjoint
  $ f^* : \CC \to \mathcal{D} $,
  its right adjoint $f_*$ is lax $\En$-monoidal.
  In particular, $f_* (\o)$ is an $\En$-ring,
  so we obtain a factorization of $f^*$ as an $\mathbb{E}_{n-1}$-monoidal left adjoint into
  \[ \CC \xrightarrow{ - \otimes f_*\o} \Mod ( \CC ; f_*\o ) \xrightarrow{g^*} \mathcal{D}. \]
  We will say that the adjunction $f$ is \emph{$0$-affine} if $g$ is an adjoint equivalence.
  Note that being $0$-affine is a property of the underlying monoidal functor. \todo{This really should have a cite to Lurie for the functor Mod out of some cartesian fibration.}
\end{cnstr}

The main result of this appendix is a convenient criterion for $f$ to be $0$-affine.
Before we can state that result we make a preparatory definition.

\begin{dfn}
  Given a monoidal left adjoint $f^* : \CC \to \mathcal{D}$,
  we may consider the projection map,
  \[ X \otimes f_*Y \xrightarrow{\eta_f} f_*f^*(X \otimes f_*Y) \simeq f_*(f^*X \otimes f^*f_*Y) \xrightarrow{\epsilon_f} f_*(f^*X \otimes Y). \]
  We will say that $f$ \emph{satisfies the projection formula at $X$},
  if for all $Y$ the projection map is an equivalence.
  If $f$ satisfies the projection formula at $X$ for all $X$, then we say that $f$ \emph{satisfies the projection formula}.
\end{dfn}

\begin{prop} \label{prop:rigid}
  Given a monoidal left adjoint $f^* : \CC \to \mathcal{D}$ between presentable categories, the following statements are true:
  \begin{enumerate}
  \item $f$ satisfies the projection formula for dualizable objects in $\CC$.
  \item If $\CC$ is rigidly generated and $\mathcal{D}$ has compact unit, then $f_*$ preserves colimits.
  \item If $\CC$ is rigidly generated  and $f_*$ preserves colimits, then $f$ satisfies the projection formula.
  \item $f_*$ is conservative if and only if the essential image of $f^*$ contains a family of generators.
  \item If $f$ satisfies the projection formula, $f_*$ preserves colimits and $f_*$ is conservative, then $f$ is $0$-affine.
  \end{enumerate}
\end{prop}

Our arguments follow those from \cite{MNN} fairly closely, though the hypotheses are somewhat different. \todo{Are there more things to cite here?} Before proceeding to the proof of this proposition we give some examples to demonstrate its effectiveness.



\begin{exm} \label{exm:fil-ctau}
In appendix \Cref{app:fil} we study the category $\Sp^{\Fil}$ of filtered spectra, where we denote the shift map $\o(-1) \to \o$ by $\tau$.
  The associated graded functor $\Gr : \Sp^{\Fil} \to \Sp^{\Gr}$ satisfies the conditions of \Cref{prop:rigid}, so we obtain an equivalence of presentably symmetric monoidal categories
  \[ \Sp^{\Gr} \cong \Mod ( \Sp^{\Fil} ; C\tau ). \]
  Here $C\tau$ acquires a commutative algebra structure from the fact that it equivalent to the image of the unit under the right adjoint of $\Gr$.
\end{exm}

\begin{exm} \label{exm:fil-inv-tau}
  Again working with $\Sp^{\Fil}$,
  the realization functor $\Re : \Sp^{\Fil} \to \Sp$ satisfies the conditions of \Cref{prop:rigid}, so we obtain an equivalence of presentably symmetric monoidal categories
  \[ \Sp \cong \Mod ( \Sp^{\Fil} ; \o[\tau^{-1}] ). \]
  Here $\o[\tau^{-1}]$ acquires a commutative algebra structure from the fact that it is equivalent to the image of the unit under $Y$ (the right adjoint of $\Re$).
\end{exm}

\begin{exm}
  Given a finite extension of characteristic zero fields $\ell/k$ we have a symmetric monoidal left adjoint $ i^* : \SH(k) \to \SH(\ell) $.
Using resolution of singularities, Nagata compactification and purity we know that both of these categories are generated by the motives of smooth projective schemes, each of which is dualizable \cite{MotivicModules}. In short,
   both the source and target categories are rigidly generated, with families of compact dualizable generators given by $\Ss^m \otimes (\P^1)^{\otimes n} \otimes X$ where $X$ is a smooth projective scheme over the base.
  
  In order to show that $i$ is 0-affine we only need to show that the image of $i^*$ contains a family of generators. Given a smooth projective $\ell$-scheme $X$ the projection formula \Cref{prop:rigid}(1) gives us an equivalence
  \[ i_*(\Ss^m \otimes (\P^1)^{\otimes n} \otimes X) \simeq \Ss^m \otimes (\P^1)^{\otimes n} \otimes i_*X. \]
  Since for field extensions $i_{\sharp} = i_*$, we learn that $i_*X = X$ where the second copy of $X$ is considered as a $k$-scheme. It will now suffice to show that $X$ is a retract of $i^*i_*X$. At the level of schemes this means looking at $ X \times_{\Spec(\ell)} \Spec(\ell) \times_{\Spec(k)} \Spec(\ell) $. Using the maps $\ell \to \ell \otimes_k \ell \to \ell$ we may conclude.

  Stated more explicitly,
  we have an equivalence of presentably symmetric monoidal categories
  \[ \SH(\ell) \cong \Mod( \SHk ; \Spec(\ell) ). \]
  Using the fact that $i$ restricts to the subcategories of Artin--Tate objects the same argument provides an equivalence of presentably symmetric monoidal categories
  \[ \SH(\ell)^{\at} \cong \Mod ( \SHkat ; \Spec(\ell) ). \]
\end{exm}

\begin{exm} \label{exm:c2-underlying}
  The category of $C_2$-spectra is rigidly generated and the underlying spectrum functor, $\Phi^e$, is essentially surjective.  We may thus apply \Cref{prop:rigid} to obtain a symmetric monoidal equivalence
  \[ \Mod( \Sp_{C_2} ; R_1) \simeq \Sp, \]
  where $R_1$ is the image of $\Ss$ under the right adjoint to $\Phi^e$.
  Since $\Phi^e$ can be described as homotopy fixed points (or homotopy orbits) composed with $(C_{2})_+ \otimes -$, we find that its right adjoint is given by tensoring with $(C_{2})_+ \simeq Ca_\sigma$. Therefore, $R_1 \simeq Ca_\sigma$.

  Similarly, with $\Phi^e$ replaced by $\Phi^{C_2}$ we obtain a symmetric monoidal equivalence
  \[ \Mod( \Sp_{C_2} ; R_2) \simeq \Sp, \]
  where $R_2$ is the image of $\Ss$ under the right adjoint to geometric fixed points.
  We can compute that $R_2 \simeq \Ss[a_\sigma^{-1}]$ using the presentation of $C_2$-spectra via an isotropy separation square. 
\end{exm}

\begin{exm} \label{exm:c2-graded}
  Taking loops on the map $i : S^1 \to B\Pic(\Sp_{C_2})$ which sends a generator to $\Ss^{\sigma}$ we get a monoidal functor $i : \Z \to \Pic(\Sp_{C_2})$. Embedding the target into $\Sp_{C_2}$ and tensoring up to spectra we obtain a monoidal left adjoint,
  \[ i^* : \Sp^{\Gr} \to \Sp_{C_2} \]
  which sends $\Ss(1)$ to $\Ss^{\sigma}$. Since $\Sp^{\Gr}$ is rigidly generated, the unit in $\Sp_{C_2}$ is compact and the representation spheres generate $\Sp_{C_2}$ we may apply \Cref{prop:rigid} to conclude that $i$ is $0$-affine. Stated more explicitly,
  we have an equivalence of presentable categories,
  \[ \Sp_{C_2} \cong \Mod( \Sp^{\Gr} ; i_*\Ss ). \]
  The graded ring $i_*\Ss$ has $n^{\mathrm{th}}$ term given by $\Map^{\Sp} ( \Ss^{n\sigma}, \Ss )$.
  Expressed in terms of stunted projective spaces we have $(i_*\Ss)_n \simeq \Sigma \R \mathrm{P}_{-\infty}^{-n-1}$. %The authors will return to the question of whether this example can be upgraded to an $\mathbb{E}_1$-deformation of spectra in the sense of the \Cref{app:def} in the future. 
\end{exm}


\begin{proof}[Proof (of \Cref{prop:rigid}(1)).]    
  Consider the following commutative diagram, which is natural in $X \in \CC^{\mathrm{dbl}}$, $Y \in \mathcal{D}$ and $Z \in \mathcal{C}$. 
  {\footnotesize
    \begin{center} \begin{tikzcd}
        \Map_{\mathcal{C}}(X^\vee \otimes Z, f_*Y) \ar[r] \ar[rr, "\simeq", bend left=10] &
        \Map_{\mathcal{D}}(f^*(X^\vee \otimes Z), f^*f_*Y) \ar[d, "\simeq"] \ar[r] &
        \Map_{\mathcal{D}}(f^*(X^\vee \otimes Z), Y) \ar[d, "\simeq"] \\        
        & \Map_{\mathcal{D}}(f^*(X^\vee) \otimes f^*Z, f^*f_*Y) \ar[d, "\simeq"] \ar[r] &
        \Map_{\mathcal{D}}(f^*(X^\vee) \otimes f^*Z, Y) \ar[d, "\simeq"] \\
        \Map_{\mathcal{C}}(Z, X \otimes f_*Y) \ar[d] \ar[uu, no head, "\simeq"] &
        \Map_{\mathcal{D}}((f^*X)^\vee \otimes f^*Z, f^*f_*Y) \ar[d, "\simeq"] \ar[r] &
        \Map_{\mathcal{D}}((f^*X)^\vee \otimes f^*Z, Y) \ar[d, "\simeq"] \\
        \Map_{\mathcal{D}}(f^*Z, f^*(X \otimes f_*Y)) \ar[r, "\simeq"] &
        \Map_{\mathcal{D}}(f^*Z , f^*X \otimes f^*f_*Y) \ar[r] &
        \Map_{\mathcal{D}}(f^*Z , f^*X \otimes Y) 
    \end{tikzcd} \end{center}
  }
  The rectangle on the left commutes due to the compatibility of dualization with the monoidal structure on $f^*$. The remaining squares commute for easier reasons. Now observe that starting in the middle of the left side and proceeding counter-clockwise gives the projection map, while proceeding clockwise given an equivalence.
\end{proof}

\begin{proof}[Proof (of \Cref{prop:rigid}(2)).]
  Since $\CC$ is rigidly generated it has a family of compact dualizable generators.
  Monoidal functors send dualizable objects to dualizable objects.
  Since the unit of $\mathcal{D}$ is compact we learn that $f^*$ sends dualizable objects to compact objects.
  Thus, since $f^*$ sends a family of compact generators to compact objects its right adjoint $f_*$ preserves colimits.
\end{proof}

\begin{proof}[Proof (of \Cref{prop:rigid}(3)).]
  The projection formula asks that the natural projection map
  \[ X \otimes f_*Y \to f_*(f^*X \otimes Y) \]
  be an equivalence.
  Using the hypotheses that $f_*$ preserves colimits and $\CC$ is rigidly generated we can reduce to the case where $X$ is a compact dualizable generator. 
  We may now use \Cref{prop:rigid}(1) to conclude.
\end{proof}

\begin{proof}[Proof (of \Cref{prop:rigid}(4)).]
  Clear.
\end{proof}

\begin{proof}[Proof (of \Cref{prop:rigid}(5)).]
  We begin by showing that $g^*$ is fully faithful.
  This is equivalent to showing that the unit map
  $X \to g_*g^*X$ is an equivalence.
  Applying the projection formula with $Y = \o$, we can conclude that this is true for induced $f_*\o$-modules.
  Since $f_*$ preserves colimits and $\Mod( \CC ; f_*\o)$ is generated by induced $f_*\o$-modules this is sufficient to conclude.

  Now, we show that $g^*$ is essentially surjective.
  By \Cref{prop:rigid}(4) we know the essential image of $f^*$ contains a family of generators (and $g^*$ has the same property).
  Using fully-faithfulness we can now conclude that $g^*$ is essentially surjective.  
\end{proof}

We close the appendix with another useful lemma.

\begin{lem} \label{lem:tensor-mod}
  Suppose that $\CC$ and $\D$ are stable presentably symmetric monoidal categories, and let $R$ be a commutative algebra in $\CC$.
  Given a symmetric monoidal left adjoint $f^* : \CC \to \mathcal{D}$, 
  there is an equivalence of presentably symmetric monoidal categories
  \[\Mod(\D; f^*R) \simeq \Mod(\CC; R) \otimes_\CC \D.\]
\end{lem}

\begin{proof}
  This follows from \cite[Theorem 4.8.5.16]{HA} after unraveling the definitions.
\end{proof}

\begin{exm}  
  Given a stable presentably symmetric monoidal category $\CC$ and a commutative algebra $R \in \CC^{\Fil}$, this lemma provides an equivalence
  \[ \Sp^{\Gr} \otimes_{\Sp^{\Fil}} \Mod ( \CC^{\Fil} ; R ) \simeq \CC^{\Gr} \otimes_{\CC^{\Fil}} \Mod ( \CC^{\Fil} ; R ) \simeq \Mod ( \CC^{\Gr} ; \Gr(R) ) \]
\end{exm}
